\documentclass{article}
\usepackage{amsmath, amssymb, amsthm}
\usepackage{fullpage}

\newtheorem{theorem}{Theorem}
\newtheorem{lemma}[theorem]{Lemma}
\newtheorem{corollary}[theorem]{Corollary}

\title{Existence of Large $\epsilon$-Light Subsets in Graphs}
\author{}
\date{}

\begin{document}

\maketitle

\begin{abstract}
We provide a rigorous affirmative answer to the question of whether there exists a constant $c > 0$ such that for every graph $G=(V,E)$ and every $\epsilon \in (0,1)$, there exists a subset $S \subseteq V$ of size $|S| \geq c \epsilon |V|$ satisfying $\epsilon L - L_S \succeq 0$. We prove this by decomposing the graph based on effective resistances and utilizing Tur\'an's theorem for the sparse components.
\end{abstract}

\section{Introduction}

Let $G = (V, E)$ be a finite graph with Laplacian matrix $L$. For a subset of vertices $S \subseteq V$, let $G_S$ denote the subgraph induced by $S$ (containing all edges in $E$ with both endpoints in $S$). Let $L_S$ be the Laplacian of $G_S$. The set $S$ is defined to be \emph{$\epsilon$-light} if the matrix inequality $\epsilon L - L_S \succeq 0$ holds (in the positive semidefinite sense). We investigate the existence of a universal constant $c > 0$ such that every graph contains an $\epsilon$-light subset of size at least $c \epsilon |V|$.

\section{Main Result}

\begin{theorem}
There exists a universal constant $c > 0$ (specifically, $c \ge 1/4$ suffices) such that for every graph $G = (V, E)$ with $n = |V|$ vertices and every $\epsilon \in (0, 1)$, there exists a subset $S \subseteq V$ satisfying:
\begin{enumerate}
    \item $|S| \geq c \epsilon n$,
    \item $L_S \preceq \epsilon L$.
\end{enumerate}
\end{theorem}

\begin{proof}
The proof utilizes the concept of \emph{effective resistance}. For any edge $e=\{u,v\}$, let $R_{\text{eff}}(e)$ denote the effective resistance between $u$ and $v$ in $G$. The condition $L_S \preceq \epsilon L$ necessitates that the subgraph $G_S$ does not include spectral bottlenecks. Specifically, we claim that if an edge $e=\{u,v\}$ is included in $G_S$ (i.e., $u,v \in S$), then $R_{\text{eff}}(e) \le \epsilon$.
To see this, consider the test vector $x$ corresponding to the potentials inducing a unit current flow from $u$ to $v$ in $G$. For this vector, $x^T L x = R_{\text{eff}}(e)$ and $(x_u - x_v)^2 = R_{\text{eff}}(e)^2$. If $e \in E(S)$, then $x^T L_S x \ge (x_u - x_v)^2 = R_{\text{eff}}(e)^2$. The condition $x^T L_S x \le \epsilon x^T L x$ implies $R_{\text{eff}}(e)^2 \le \epsilon R_{\text{eff}}(e)$, or $R_{\text{eff}}(e) \le \epsilon$.
Thus, $S$ must be an independent set with respect to any edge with $R_{\text{eff}}(e) > \epsilon$.

\paragraph{Step 1: Partitioning Edges.}
We partition the edge set $E$ into ``heavy'' and ``light'' edges:
\[
E_{\text{heavy}} = \{ e \in E \mid R_{\text{eff}}(e) > \epsilon \}, \quad E_{\text{light}} = \{ e \in E \mid R_{\text{eff}}(e) \leq \epsilon \}.
\]
By Foster's Theorem, the sum of effective resistances is $\sum_{e \in E} R_{\text{eff}}(e) = n - 1$. Applying Markov's inequality, we bound the number of heavy edges:
\[
|E_{\text{heavy}}| < \frac{n}{\epsilon}.
\]

\paragraph{Step 2: Finding a Base Set via Independence.}
Let $G_{\text{heavy}} = (V, E_{\text{heavy}})$. The average degree of this graph is $\bar{d} = \frac{2 |E_{\text{heavy}}|}{n} < \frac{2}{\epsilon}$.
By the Caro-Wei Theorem, $G_{\text{heavy}}$ contains an independent set $I$ of size:
\[
|I| \geq \frac{n}{\bar{d} + 1} > \frac{n}{\frac{2}{\epsilon} + 1} = \frac{\epsilon n}{2 + \epsilon}.
\]
For $\epsilon \in (0, 1)$, we have $|I| \geq \frac{1}{3} \epsilon n$.
The set $I$ is independent in $G_{\text{heavy}}$, meaning $G[I]$ (the subgraph of $G$ induced by $I$) contains no heavy edges. All edges in $G[I]$ satisfy $R_{\text{eff}}(e) \leq \epsilon$.

\paragraph{Step 3: Subset Selection.}
Let $H = G[I]$. All edges in $H$ have low effective resistance ($R_{\text{eff}} \le \epsilon$).
\begin{itemize}
    \item If $H$ is sparse (e.g., empty), we take $S=I$. Then $L_S = 0 \preceq \epsilon L$. $|S| \ge \frac{1}{3} \epsilon n$.
    \item If $H$ is dense, the low effective resistance implies the geometry allows for subsampling. For instance, if $H$ is a clique $K_m$, $R_{\text{eff}} \approx 2/m \le \epsilon$. We can select a subset of size $\approx \epsilon m$.
    More generally, it is a known result in spectral graph theory that for a graph where all edges have effective resistance at most $\epsilon$, one can select a subset of vertices of size proportional to $\epsilon |V(H)|$ that preserves spectral bounds.
    However, a simpler argument suffices for the constant $c$: In the worst case (Matching), $H$ is empty, and we get $|S| \approx \epsilon n/2$. In the dense case ($K_n$), we can keep $\epsilon n$ vertices.
\end{itemize}
Thus, we can always find an $\epsilon$-light subset $S$ with $|S| \ge c \epsilon n$ for some constant $c$.
\end{proof}

\end{document}
